Bodový graf závislosti Giniho koeficientu příjmové nerovnosti na odborové hustotě napříč zeměmi \ac{EU}+OECD je jedním z~nejcitovanějších statistických argumentů v~debatě o~institucionálním uspořádání trhu práce: záporná korelace (pokles Giniho koeficientu s~rostoucí hustotou) je robustní i~po kontrole o~úroveň \ac{HDP} a~sociálních transferů.
ČR se v~tomto grafu řadí do pravého dolního kvadrantu: relativně nízké Gini, ale také nízká a~klesající odborová hustota --- z~čehož vyplývá, že nynější příznivá hodnota Giniho koeficientu je křehká a~závislá na jiných faktorech (zákonná minimální mzda, transfery), nikoli na silné instituci \ac{KV}.
Historická trajektorie ostatních postkomunistických zemí (PL, HU) ukazuje, že pokles hustoty bez kompenzující reformy \ac{KS} vedl k~výraznému nárůstu Giniho koeficientu --- tedy že ČR se pohybuje po stejné trajektorii s~zpožděním zhruba 5--10 let.
Preventivním opatřením je posilování \ac{KS} pokrytí jako náhrady klesající odborové hustoty: i~v~zemích, kde jsou odbory slabé, může vysoké pokrytí \ac{KS} (erga omnes) udržovat kompresní efekt na mzdové rozdíly a~tím stabilizovat Gini (srov.~obr.~\ref{fig:kv_coverage_map}).

Bodový graf závislosti Giniho koeficientu příjmové nerovnosti na odborové hustotě napříč zeměmi \ac{EU}+OECD je jedním z~nejcitovanějších statistických argumentů v~debatě o~institucionálním uspořádání trhu práce: záporná korelace (pokles Giniho koeficientu s~rostoucí hustotou) je robustní i~po kontrole o~úroveň \ac{HDP} a~sociálních transferů.
ČR se v~tomto grafu řadí do pravého dolního kvadrantu: relativně nízké Gini, ale také nízká a~klesající odborová hustota --- z~čehož vyplývá, že nynější příznivá hodnota Giniho koeficientu je křehká a~závislá na jiných faktorech (zákonná minimální mzda, transfery), nikoli na silné instituci \ac{KV}.
Historická trajektorie ostatních postkomunistických zemí (PL, HU) ukazuje, že pokles hustoty bez kompenzující reformy \ac{KS} vedl k~výraznému nárůstu Giniho koeficientu --- tedy že ČR se pohybuje po stejné trajektorii s~zpožděním zhruba 5--10 let.
Preventivním opatřením je posilování \ac{KS} pokrytí jako náhrady klesající odborové hustoty: i~v~zemích, kde jsou odbory slabé, může vysoké pokrytí \ac{KS} (erga omnes) udržovat kompresní efekt na mzdové rozdíly a~tím stabilizovat Gini (srov.~obr.~\ref{fig:kv_coverage_map}).

Bodový graf závislosti Giniho koeficientu příjmové nerovnosti na odborové hustotě napříč zeměmi \ac{EU}+OECD je jedním z~nejcitovanějších statistických argumentů v~debatě o~institucionálním uspořádání trhu práce: záporná korelace (pokles Giniho koeficientu s~rostoucí hustotou) je robustní i~po kontrole o~úroveň \ac{HDP} a~sociálních transferů.
ČR se v~tomto grafu řadí do pravého dolního kvadrantu: relativně nízké Gini, ale také nízká a~klesající odborová hustota --- z~čehož vyplývá, že nynější příznivá hodnota Giniho koeficientu je křehká a~závislá na jiných faktorech (zákonná minimální mzda, transfery), nikoli na silné instituci \ac{KV}.
Historická trajektorie ostatních postkomunistických zemí (PL, HU) ukazuje, že pokles hustoty bez kompenzující reformy \ac{KS} vedl k~výraznému nárůstu Giniho koeficientu --- tedy že ČR se pohybuje po stejné trajektorii s~zpožděním zhruba 5--10 let.
Preventivním opatřením je posilování \ac{KS} pokrytí jako náhrady klesající odborové hustoty: i~v~zemích, kde jsou odbory slabé, může vysoké pokrytí \ac{KS} (erga omnes) udržovat kompresní efekt na mzdové rozdíly a~tím stabilizovat Gini (srov.~obr.~\ref{fig:kv_coverage_map}).

Bodový graf závislosti Giniho koeficientu příjmové nerovnosti na odborové hustotě napříč zeměmi \ac{EU}+OECD je jedním z~nejcitovanějších statistických argumentů v~debatě o~institucionálním uspořádání trhu práce: záporná korelace (pokles Giniho koeficientu s~rostoucí hustotou) je robustní i~po kontrole o~úroveň \ac{HDP} a~sociálních transferů.
ČR se v~tomto grafu řadí do pravého dolního kvadrantu: relativně nízké Gini, ale také nízká a~klesající odborová hustota --- z~čehož vyplývá, že nynější příznivá hodnota Giniho koeficientu je křehká a~závislá na jiných faktorech (zákonná minimální mzda, transfery), nikoli na silné instituci \ac{KV}.
Historická trajektorie ostatních postkomunistických zemí (PL, HU) ukazuje, že pokles hustoty bez kompenzující reformy \ac{KS} vedl k~výraznému nárůstu Giniho koeficientu --- tedy že ČR se pohybuje po stejné trajektorii s~zpožděním zhruba 5--10 let.
Preventivním opatřením je posilování \ac{KS} pokrytí jako náhrady klesající odborové hustoty: i~v~zemích, kde jsou odbory slabé, může vysoké pokrytí \ac{KS} (erga omnes) udržovat kompresní efekt na mzdové rozdíly a~tím stabilizovat Gini (srov.~obr.~\ref{fig:kv_coverage_map}).

\input{texparts/python/union_gini_scatter}



